% intro.tex

\chapter{绪论}
\echapter{Introduction}

\section{研究背景与意义}
\esection{Research Background and Significance}

随着物联网、工业互联网、智慧交通和能源管理系统的发展，实际生产和生活中持续产生大量时间序列数据。这类数据通常按照固定或近似固定的时间间隔采集，例如电力系统中的负荷变化、交通系统中的道路占用率、气象系统中的温度和湿度变化等。对这些历史数据进行建模，并预测未来一段时间内的变化趋势，是时序数据分析中的重要任务。时序数据预测不仅关系到预测精度，也关系到实际系统中的调度决策、风险预警和资源配置。

与单步预测相比，多步预测更容易受到误差累积、周期变化和分布漂移的影响。在真实场景中，时间序列往往不是简单平稳的曲线，而是同时包含趋势项、周期项、局部波动以及噪声。以交通流量为例，工作日早晚高峰会形成明显周期，但突发拥堵又会带来短时间内的剧烈变化；以气象数据为例，不同气象变量之间存在一定物理联系，但这些联系在不同季节和时间段中也会发生变化。因此，一个有效的时序数据预测模型需要同时关注局部变化和整体趋势，并尽可能利用多变量之间的相关信息。

近年来，深度学习方法在时间序列预测任务中得到广泛应用。循环神经网络、卷积神经网络、Transformer 以及 MLP 类模型分别从不同角度改进了序列建模能力。其中，Transformer 类模型依靠注意力机制建模较大范围内的时间依赖，但在输入序列较长时通常计算成本较高；线性模型和 MLP 类模型结构较轻，在效率方面具有优势，但需要合理的结构设计才能充分提取复杂时间模式。如何在预测精度、模型复杂度和训练效率之间取得平衡，是时序数据预测模型研究中的一个重要问题。

本文以 TimeMixer 模型为研究基础。TimeMixer 是一种轻量级多尺度时序预测模型，通过多尺度序列构建、趋势季节分解和跨尺度信息混合实现预测\cite{wang2023timemixer}。该模型整体结构清晰，计算效率较高，适合作为本科毕业设计中进行结构改进和实验分析的基础模型。通过阅读代码和实验记录可以发现，原始 TimeMixer 在使用过程中仍有进一步改进空间：一是下采样阶段主要依赖平均池化，可能平滑掉局部突变信息；二是趋势分解使用固定窗口，难以适应不同数据集中的多周期模式；三是默认通道独立策略对变量间关系利用不足。针对这些问题，本文设计并实现了三个改进模块，并通过完整消融实验分析其作用。

本课题的意义主要体现在两个方面。理论方面，本文围绕轻量级多尺度预测模型，探索可学习下采样、多核门控分解和通道混合在时序数据预测中的作用，为理解 TimeMixer 这类模型的结构改进提供实验参考。应用方面，电力、交通和气象等场景都需要较稳定的预测模型，本文的实验分析可为类似场景下选择和改进轻量级预测模型提供一定借鉴。需要说明的是，本文重点不在于声称模型全面超过已有先进方法，而在于围绕 TimeMixer 进行多尺度混合增强设计，并通过消融实验如实分析各模块的贡献和不足。

\section{国内外研究现状}
\esection{Related Work}

\subsection{基于统计模型的时间序列预测方法}
\esubsection{Statistical Forecasting Methods}

时间序列预测最早主要依赖统计模型。常见方法包括自回归模型、移动平均模型、自回归滑动平均模型以及差分自回归滑动平均模型等。这类方法通常假设序列具有较强的线性关系或经过差分后接近平稳，在许多传统业务场景中具有较好的可解释性。统计模型的优点是结构简单、参数含义明确、训练和推理成本较低，因此在样本规模较小或预测对象较单一时仍然有一定应用价值。

但是，传统统计模型也存在明显限制。首先，它们通常需要人工判断平稳性、季节性和差分阶数，对使用者的经验依赖较强。其次，对于高维多变量序列，传统模型难以充分刻画变量之间复杂的非线性关系。最后，实际数据常常存在突变、噪声和多重周期，单一统计假设难以覆盖所有情况。因此，在公开时序预测数据集和复杂工业场景中，统计模型更多被用作基础方法或理论参照。

\subsection{基于深度学习的时间序列预测方法}
\esubsection{Deep Learning Forecasting Methods}

深度学习的发展推动了时间序列预测方法的变化。循环神经网络能够按照时间顺序处理序列数据，长短期记忆网络通过门控结构缓解了普通循环网络的梯度消失问题\cite{hochreiter1997lstm}，在较早的序列建模任务中应用较多。卷积神经网络则通过局部卷积核提取时间邻域特征，具有较好的并行计算能力。后来，Transformer 结构被引入时间序列预测任务，通过自注意力机制建模长距离依赖，使模型能够在较长输入窗口中选择重要时间位置。

针对多步时序预测中输入窗口和预测步长较大、注意力计算成本较高的问题，研究者提出了多种改进结构。例如 Informer 引入稀疏注意力机制以降低复杂度\cite{zhou2021informer}；Autoformer 将分解思想引入深度预测模型，用自相关机制提取周期信息\cite{wu2021autoformer}；FEDformer 从频域角度改进周期建模能力\cite{zhou2022fedformer}。这些方法说明，时序数据预测不仅需要一般的序列建模能力，还需要结合时间序列本身的趋势、季节性和周期结构。

\subsection{基于分解与多尺度混合的预测方法}
\esubsection{Decomposition and Multi-scale Forecasting Methods}

近年来，一些研究开始重新关注简单结构在时序数据预测中的作用。DLinear 等方法表明，在部分数据集上，分解加线性映射的结构也能取得有竞争力的效果\cite{zeng2023transformers}。这说明复杂模型并不总是必要，合理利用时间序列的趋势和季节模式同样重要。PatchTST 等方法将相邻时间点划分为局部片段，从而兼顾局部模式和较长范围依赖\cite{nie2023time}；MLP 类方法则尝试用更轻量的结构实现时间维度和变量维度的信息混合\cite{chen2023tsmixer}。

TimeMixer 属于这一类轻量级多尺度建模方法。它通过下采样构建不同时间分辨率的序列，再在过去可分解混合模块中对不同尺度的趋势项和季节项进行交互，最后融合多尺度预测结果。与单一尺度模型相比，多尺度结构能够同时关注局部波动和整体趋势；与复杂注意力模型相比，TimeMixer 的计算结构相对简洁。但是，正因为模型追求轻量化，一些设计也较为朴素，例如平均池化下采样、固定窗口分解以及通道独立处理。这些特点为本文的改进提供了切入点。

\section{TimeMixer模型存在的问题}
\esection{Limitations of TimeMixer}

结合原始代码和本文实验记录，本文认为原始 TimeMixer 主要存在以下三方面不足。

第一，下采样方式较为固定。TimeMixer 在构建多尺度输入时默认采用平均池化。平均池化计算简单且没有额外参数，但它对窗口内时间点赋予相同权重，本质上是一种平滑操作。当序列中存在尖峰、骤降或局部突变时，平均池化可能削弱这些高频信息。对于交通流量、电力负荷等具有突发变化的数据，这种信息损失可能影响后续预测。

第二，分解窗口缺乏自适应性。原始模型使用固定大小的滑动平均窗口分离趋势项和季节项。固定窗口能够提供稳定的趋势估计，但实际序列往往同时包含短期波动、中期周期和长期变化。单一窗口难以同时适应这些不同尺度的模式，在周期结构复杂的数据集上可能限制模型表达能力。

第三，通道独立策略对变量关系利用不足。TimeMixer 为了提高计算效率，通常将不同变量拆分为独立序列进行建模。这种方式能够降低多变量建模难度，但也可能忽略变量之间的相关性。例如气象数据中的温度、湿度、风速之间存在联系，交通数据中不同传感器之间也可能存在空间关联。完全独立建模可能使模型难以利用这部分信息。

\section{本文主要研究内容与章节安排}
\esection{Main Work and Organization}

针对上述问题，本文围绕 TimeMixer 进行了三项结构改进。

第一，设计深度可分离卷积下采样模块。该模块使用可学习的一维卷积替代平均池化，使下采样过程从固定平滑变为可学习特征提取。为控制参数量，模块采用逐通道卷积和逐点卷积组合，并加入 LayerNorm 和 GELU 激活。

第二，设计多核门控分解模块。该模块并行使用多个滑动平均窗口提取不同尺度的趋势估计，再通过门控网络为不同窗口分配融合权重。当前代码中实际窗口大小为 $(3, 11, 37)$，论文后续均以该设置为准。

第三，引入通道混合模块。该模块在输入归一化后、通道独立拆分前，使用低秩 MLP 对变量维度进行一次残差映射，从而在保持轻量化结构的同时引入变量间信息交互。

为分析不同模块的独立作用和组合效果，本文构造 M0、M1、M2、M3、M12、M13、M23 和 M4 共 8 个模型变体，在 8 个公开数据集和 4 种预测长度上进行实验。实验结果表明，多核门控分解模块是本文最稳定的有效改进；深度可分离下采样和通道混合在部分场景中有一定作用，但稳定性不足；完整模型 M4 结构最完整，但并不总是取得最优结果。

本文后续章节安排如下：第 2 章介绍时序数据预测任务、时间序列分解、多尺度建模思想以及原始 TimeMixer 模型结构；第 3 章详细介绍本文提出的三个改进模块和消融实验变体设计；第 4 章说明实验设置、数据集、评价指标和实验结果，并对结果进行分析；第 5 章总结全文工作，说明不足并给出后续改进方向。
